%% kapa: exam=1 round=22 year=2011 period=1 form=A subject=civil_law subject_ko=민법 seq=38 number=38
%% source: Q-Net 감정평가사 기출문제 (cst00309) artlSeq=5036790 file=extracted/1008/5036790/2011년 제22회 감정평가사 1차시험 기출문제/1교시.hwp
%% conversion: kordoc 4.24.51 -> tools/gen1.py (auto); review=proofread-sonnet; flags=none
\begin{question}{38}
물권의 소멸사유로서 혼동에 관한 설명으로 옳은 것은?
\begin{choices}
  \item 乙이 甲소유의 토지에 지상권을 취득한 다음 그 지상권을 목적으로 丙에게 저당권을 설정한 후 甲으로부터 그 토지를 매수한 경우, 乙의 지상권은 소멸한다.
  \item 물상보증인 丁소유의 토지에 채무자 甲의 채권자 乙이 1순위, 채권자 丙이 2순위의 저당권자인 경우, 乙이 단독으로 丁을 상속한 때에는 乙의 저당권은 소멸한다.
  \item 甲소유의 토지를 점유하고 있는 乙이 甲으로부터 그 토지를 매수한 경우, 乙의 점유권은 소멸한다.
  \item 채무자 甲소유의 토지에 채권자 乙이 1순위, 채권자 丙이 2순위의 저당권자인 경우, 丙이 甲으로부터 그 토지 소유권을 취득한 때에는 丙의 저당권은 소멸하지 않는다.
  \item 甲소유의 부동산에 乙이 전세권을 취득한 후 丙이 근저당권을 취득한 경우, 乙이 그 부동산을 증여받았더라도 乙의 전세권은 소멸하지 않는다.
\end{choices}
\end{question}
\answer{5}
